\section{第三日第一讲}

我们还没有谈过球簇，它在 BZSV 的故事里有重要地位。这里我们不打算一上来就给出定义，而是通过考察「什么样的簇让我们能够推广此前见到的调和分析」来引出它。

设 $G$ 是非 archimedean 域 $F$ 上的分裂约化群，$O$ 是赋值环，剩余域为 $\mathbb F_q$。记
\[
  H_G:=H(G(F),G(O))
\]
为 Hecke 代数，把它看作 $G(F)$ 上紧支集 $G(O)$-双不变函数的空间，乘法由 $G$ 上 Haar 测度的卷积给出。由 Satake 同构，有同构
\[
  H_G\xrightarrow{\ \sim\ }\mathbb C[\Rep\check G],\qquad f\mapsto\hat f,
\]
其中右端指张量范畴 $\Rep_{\mathbb C}\check G$ 的 Grothendieck 环 $K_0$ 上的群代数。

现在作升级。设 $X$ 是带 $G$ 作用的光滑仿射空间。可以考虑紧支集局部常值函数空间 $C_c^\infty(X(F))$，我们称之为 $X$ 的「Schwartz 空间」，记作 $\mathcal S(X(F))$。它自然是 $G(F)$-表示。再考虑右 $G(O)$-不变子空间
\[
  \mathcal S(X(F))^{G(O)}:=C_c^\infty(X(F))^{G(O)},
\]
它是 Hecke 代数 $H_G$ 右作用的模。（这叫 $X(F)$ 的「球」Schwartz 空间。）但我们此前也看到，其实应当把这个向量空间看作内积空间：
\[
  \mathcal S(X(F))^{G(O)}\subset\big(L^2(X(F)),\langle\ ,\ \rangle\big),
\]
其中乘积按 $X(F)$ 上的不变测度积分。（另一种做法是把 $L^2$ 空间取成 $X(F)$ 上的半密度，此时球 Hecke 代数中的函数 $f$ 要写成 $f|dx|^{1/2}$。）记
\[
  \mathbf 1_X:=\mathbf 1_{X(O)}
\]
为 $X(O)$ 的示性函数。这是 Schwartz 空间中的「基本向量」。（一般地，这个基本函数的公式在 $X$ 光滑时是好的。对于非光滑簇（不在本文范围内），基本向量应当与 $X$ 的环路空间的 IC 层相关联。）

在此前的讲座与练习中，我们已经遇到过球 Schwartz 空间中内积的两个例子：

\begin{exmp}
  设 $X=\mathbb A^1$，$G=\mathbb G_m$ 按伸缩作用。$\mathbb A^1$ 中的 $G(O)$-轨道恰是环带 $\{x\in\mathbb A^1(F):|x|=q^{-n}\}_{n\in\mathbb Z}$，而球 Schwartz 空间由这些环带的示性函数的有限组合构成。Hecke 代数是 $H_G\simeq\mathbb C[\mathbb Z]\xrightarrow{\sim}\mathbb C[\mathbb C^\times]$。（这里不可避免地有个记号冲突：$\mathbb C[\mathbb Z]$ 指 $\mathbb Z$ 的群代数，而 $\mathbb C[\mathbb C^\times]$ 指 $\mathbb C$ 上 $\mathbb C^\times=\mathbb G_m$ 的代数函数。）由此可见 Satake 变换是一种 Fourier 变换，我们把它记作 $h\in H_G\mapsto\hat h\in\mathbb C[\mathbb C^\times]$。于是 $\mathcal S(X(F))$ 中的内积用 Satake 变换写出来如下：
  \begin{equation}
    \langle h\star\mathbf 1_X,\mathbf 1_X\rangle
    =\int_{S^1\subset\mathbb C^\times}\frac{\hat h(z)}{(1-q^{-1/2}z)(1-q^{-1/2}z^{-1})}\,d^\times z
    \label{eq:8-2}
  \end{equation}
  其中 $d^\times z=\frac{dz}{z}$ 是乘法 Haar 测度。注意我们做作业时已经见过这个 q-特征标。
\end{exmp}

\begin{defn}
  量 $\dfrac{d^\times z}{(1-q^{-1/2}z)(1-q^{-1/2}z^{-1})}$ 称为 $\mathbf 1_{\mathbb A^1}$ 的 \textbf{Plancherel 密度}。
\end{defn}

让我们回顾这个 q-特征标在「对偶」（即 $B$）侧是如何出现的。一般地 $M=T^*X$，所以这里 $M=T^*\mathbb A^1$。对偶地，我们得到 $\check G=\mathbb G_m$，它作用在 $\check M$ 上，而 $\check M$ 恰好也等于 $T^*\mathbb A^1$。作用按权 $(1,-1)$ 进行，即
\[
  \lambda.(x,y)\mapsto(\lambda x,\lambda^{-1}y)\in T^*\mathbb A^1=\mathbb A^1\times(\mathbb A^1)^*.
\]
（这里 $\mathbb G_m$ 与 $T^*\mathbb A^1$ 各自带有一个例外的自对偶；一般地，$\check G$ 与 $\check M$ 会和 $G$ 与 $M$ 很不一样。）

回顾 q-特征标的定义：
\begin{equation}
  \mathrm{q\text{-}tr}(z)=\sum_{i\ge0}q^{-i/2}\operatorname{tr}\big(z\,|\,\mathbb C[\check M]_i\big)
  =\frac{1}{\det(I-q^{-1/2}z\,|\,\check M)},
  \label{eq:8-4}
\end{equation}
其中 $\mathbb C[\check M]_i=\Sym^i(\check M^\vee)=\Sym^i(\check M)$（用到 $\check M=\mathbb A^1\times(\mathbb A^1)^*$ 作为向量空间的线性自对偶）。这里的分次来自我们即将称之为 $G_{\mathrm{gr}}$ 的作用，它把 $\check M$ 的两个 $\mathbb A^1$ 因子按同一比例伸缩（即 $G_{\mathrm{gr}}$ 按权 $(1,1)$ 作用）。由于 $G_{\mathrm{gr}}$ 同构于 $\mathbb G_m$，这可能引起混淆；所以我们在此强调：$\check M$ 上有\textbf{两个不同}的 $\mathbb G_m$ 作用，即 $\check G=\mathbb G_m$ 的作用与 $G_{\mathrm{gr}}=\mathbb G_m$ 的作用。

总结一下：表达式 $\dfrac{1}{(1-q^{-1/2}z)(1-q^{-1/2}z^{-1})}$ 是 $\check M=T^*\mathbb A^1$ 在 $G_{\mathrm{gr}}$-分次下的 q-特征标。我们把它记作 $L(z,\check M)$。一般地，Plancherel 密度形如
\[
  \text{Plancherel 密度}=(\text{辛簇}\ \check M\ \text{的 q-特征标})\,d^\times z=L(z,\check M)\,d^\times z,
\]
于是 (8.2) 可以重写为
\begin{equation}
  \langle h\star\mathbf 1_X,\mathbf 1_X\rangle=\int_{S^1\subset\mathbb C^\times}\hat h(z)\,L(z,\check M)\,d^\times z.
  \label{eq:8-5}
\end{equation}
我们希望这条公式对更一般的一类 $G$-簇 $X$ 也成立，其中 $M=T^*X$，对偶为 $\check G_X,\check M$。（注意 Plancherel 密度的公式里只出现 $\check M$——特别地，若多个 $G$-簇 $X$ 给出同一个 $\check G_X$ 与 $\check M$，则它们基本向量的 Plancherel 密度应当相同。）

\textbf{注}（给算术背景的读者）。我们的公式里有很多 $q^{1/2}$；这提示在 $\mathbb C$ 以外的域上工作时，必须选定一个 $q$ 的平方根。算术设定中的这种歧义由所谓的 $C$-群（而不是 $L$-群）来消解。当然在 $\mathbb C$ 上，可以取正的 $q^{1/2}$。

\begin{exmp}
  设 $X=H$，$G=H\times H$，其中 $H$ 是分裂约化群。双模结构为 $(h_1,h_2)\in G:x\in X\mapsto h_1^{-1}xh_2$。Schwartz 空间恰好就是 $H$ 的 Hecke 代数：
  \[
    \mathcal S(X)^{G(O)}\cong H_H.
  \]
  在对偶侧，我们得到 $\check G=\check H\times\check H$ 与 $\check G_X=\check H$。于是有反对角嵌入
  \[
    \check G_X=\check H\xrightarrow{(C,\mathrm{Id})}\check H\times\check H=\check G,\qquad a\mapsto(a^{-1},a),
  \]
  其中 $C$ 表示 Cartan 对合。于是在经典 Satake 变换下，Hecke 代数的等同变为
  \[
    H_G=H_H\otimes H_H\cong\mathbb C[\check H]^{\check H},\qquad h_1\otimes h_2\mapsto\hat h_1(z^{-1})\cdot\hat h_2(z).
  \]
  于是此情形下的内积读作
  \begin{equation}
    \langle(h_1\otimes h_2)\star\mathbf 1_H,\mathbf 1_H\rangle
    =\int_{\check H_c}\frac{\hat h_1(z^{-1})\hat h_2(z)}{\det(I-q^{-1}z\,|\,\check{\mathfrak h})}\,d_{Haar}(z)
    \label{eq:8-7}
  \end{equation}
  \begin{equation}
    \phantom{\langle(h_1\otimes h_2)\star\mathbf 1_H,\mathbf 1_H\rangle}
    =\int_{\check H_c}\hat h_1(z^{-1})\hat h(z)\,L(z,\check{\mathfrak h})\,d_{Haar}(z).
    \label{eq:8-8}
  \end{equation}
  下面澄清积分中出现的各项。我们在对偶群 $\check H$ 的紧形式 $\check H_c$ 上积分，$d_{Haar}(z)$ 表示 Haar 测度。记 $\check{\mathfrak h}=\mathrm{Lie}(\check H)$，$L(z,\check{\mathfrak h}):=\dfrac{1}{\det(\mathrm{Id}-q^{-1}z)}$，它表示 $\check H$ 在 $\check{\mathfrak h}$ 上伴随作用的 q-特征标。最后，$\check{\mathfrak h}$ 上的分次 $G_{\mathrm{gr}}$-作用为 $\lambda.z\mapsto\lambda^2z$，即通常伸缩作用的平方（这就是分母中出现 $q^{-1}$ 而不是 $q^{-1/2}$ 的原因）。注意积分所在的群不是 $\check G_c$ 而是 $\check H_c$，这反映了此例中 $\check G_X=\check H$ 而非 $\check G$。
\end{exmp}

如我们已提到的，有一个一般的预期：对足够好的一类 $(G,X)$（$G\circlearrowleft X$），类似的故事成立。给定 $X$，我们预期存在对偶群 $\check G_X\subset\check G$，具有如下性质。我们有作用 $\mathcal S(X(F))^{G(O)}\circlearrowleft H_G$，它给出映射
\begin{equation}
  H_G\to\End\Big(\mathcal S(X(F))^{G(O)}\Big).
  \label{eq:8-9}
\end{equation}
而在 Satake 同构 $H_G\simeq\mathbb C[\Rep\check G]$ 下，其中 $\mathbb C[\Rep\check G_X]$ 表示 Grothendieck 环 $K_0(\Rep\check G_X)$ 的群代数。$\check G$ 的表示自然限制为任意子群 $\check G'\subset\check G$ 的表示；特别地，这给出自然的环同态 $\mathbb C[\Rep\check G]\xrightarrow{\text{限制}}\mathbb C[\Rep\check G']$。我们预期 $\check G_X$ 是 $\check G$ 的约化子群，使得 $H_G$ 的作用分解为：
\begin{equation}
  \begin{array}{ccc}
    H_G & \longrightarrow & \End_{\mathbb C}\big(\mathcal S(X(F))^{G(O)}\big)\\[3pt]
    \overset{\text{Satake}}{\simeq}\Big\downarrow & & \Big\uparrow\\[3pt]
    \mathbb C[\Rep\check G] & \xrightarrow{\ \text{限制}\ } & \mathbb C[\Rep\check G_X].
  \end{array}
  \label{eq:8-10}
\end{equation}
换句话说，我们预期 Hecke 代数 $H_G\simeq\mathbb C[\Rep\check G]$ 在球 Schwartz 空间上的作用「经过」它在 $\mathbb C[\Rep\check G_X]$ 上的限制分解。（注意这个限制映射未必满射。）

我们还预期「基本向量」$\mathbf 1_{X(O)}$ 作为 $H_G$-模生成整个球 Schwartz 空间 $\mathcal S(X(F))^{G(O)}$。事实上我们预期得更多。记 $\mathbb C[\Rep\check G_X]:=H_G^X$。我们预期模 $\mathcal S(X(F))^{G(O)}$ 在 $H_G^X$ 上秩 $1$ 自由，且有指定的生成元 $\mathbf 1_{X(O)}$。换句话说，可以把 (8.10) 的右侧竖直映射看作给出典范同构
\begin{equation}
  H_G^X\xrightarrow{\ \sim\ }\End_{H_G}\Big(\mathcal S(X(F))^{G(O)}\Big),
  \label{eq:8-11}
\end{equation}
或者，由于 $\mathcal S(X(F))^{G(O)}$ 中有指定的向量，我们得到向量空间的自然同构
\begin{equation}
  H_G^X\xrightarrow{\ \sim\ }\mathcal S(X(F))^{G(O)},
  \label{eq:8-12}
\end{equation}
它把
\begin{equation}
  \mathbb C[\Rep\check G_X]\ni\mathrm{Triv}\mapsto\mathbf 1_{X(O)}\in\mathcal S(X(F))^{G(O)}.
  \label{eq:8-13}
\end{equation}
在此等同下，$\mathcal S(X(F))^{G(O)}$ 从左端继承的乘法来自融合（fusion）。（事实上，当 $X$ 不是群时，$X$ 上的卷积根本没有意义。）

最后我们指出：为代数 $\End_{H_G}(\mathcal S(X(F))^{G(O)})$（或在上述等同下，$X(F)$ 上右 $G(O)$-等变反常层的范畴）发展 Tannakian 形式化，将把我们引向 $\check G_X$ 的 Gaitsgory--Nadler 构造。

对 $h\in H_G$，记 $\hat h$ 为其经典 Satake 变换，我们把它（通过迹）看作 $\check G$ 上的代数（类）函数。我们的预期是：
\begin{equation}
  \langle\mathbf 1_{X(O)}\star h,\mathbf 1_{X(O)}\rangle
  =\int_{\check G_{X,c}\subset\check G_X}\hat h(z)\,L(z,V_X)\,d_{Haar}z,
  \label{eq:8-14}
\end{equation}
更一般地，
\begin{equation}
  \langle\mathbf 1_{X(O)}\star h_1,\mathbf 1_{X(O)}\star h_2\rangle
  =\int_{\check G_{X,c}\subset\check G_X}\hat h_1(z)\hat h_2(z)\,L(z,V_X)\,d_{Haar}z,
  \label{eq:8-15}
\end{equation}
其中 $\check G_{X,c}$ 表示 $\check G_X\subset\check G$ 的紧形式的实点，$V_X$ 表示 $\check G_X$ 的一个 $G_{\mathrm{gr}}$-分次表示，$L(z,V_X)$ 是推广了 q-特征标的适当局部 L-因子。注意 (8.15) 的右端只看到 $\hat h$ 在 $\check G_X$ 上的限制值（事实上是在 $\check G_X$ 的紧形式上的限制值），而按定义 $\hat h$ 是 $\check G$ 上的函数。

在 Satake 变换下要保持内积，就必须纳入 L-因子（而它又依赖于 $G_{\mathrm{gr}}$-分次）；这反映了导出 Satake 等价的几何本性，将在其他讲座中讨论。关键想法是：范畴化之后，内积变成 Hom（或 Ext）群，而 $G_{\mathrm{gr}}$-分次对应于 $X(F)$ 上等变层导出范畴中的一次 \textbf{shearing}（斜框）。

但我们还是回到这个问题：哪些 $G$-簇可能满足上述种种期望？结果表明，这些期望极大地限制 $X$ 的类别。特别地，它们蕴含：对一切特征标 $\chi:H_G\to\mathbb C$（等价地，对一切 $\chi\in\check A//W_G(\mathbb C)$，其中 $\check A$ 是对偶环面），
\begin{equation}
  \dim\Hom_{H_G}\big(\Func(X(F),\mathbb C),\mathbb C_\chi\big)<\infty.
  \label{eq:8-16}
\end{equation}
例如，取 $X=G$，且只考虑 $G$ 在自身上的左作用，则由 Peter--Weyl 有
\begin{equation}
  \Fun(G(F))=\int^{\oplus}\pi\otimes e_\pi\,d\pi,
  \label{eq:8-17}
\end{equation}
它对每个特征标都不满足 (8.16)（因为对多数不可约表示 $\pi$ 有 $\dim(\pi)=\infty$）。

\textbf{断言}。设 $G$ 分裂。则有限性条件 (8.16) 成立，当且仅当 $X$ 是球簇。

这就是球簇这一类进入 BZSV 框架的缘由。

\begin{defn}
  正规不可约 $G$-簇 $X$ 称为\textbf{球簇}，若某个（等价地，任一）Borel 子群 $B\subset G$ 的作用有开稠轨道。
\end{defn}

不失一般性，此后我们假定所有球簇都是拟仿射的。这不损失一般性，因为总可以化归到这种情形——代价是引入一个交换的 $G_{\mathrm m}^N$ 作用在 $X$ 上。事实上，设给定射影嵌入 $X\to\mathbb P^N$。考虑通常的满射 $V^*:=V\setminus\{0\}\to\mathbb P^N$（$V=\mathbb A^{n+1}$），并用笛卡尔图中的 $\tilde X$ 替换 $X$：
\[
  \begin{array}{ccc}
    G\times G_{\mathrm m}^N\ \circlearrowleft & \tilde X & \longrightarrow\ V^*\\
    & \Big\downarrow & \Big\downarrow\\
    G\ \circlearrowleft & X & \longrightarrow\ \mathbb P^N.
  \end{array}
\]
例如，可以把旗簇 $G/B$（带左 $G$ 作用）换成基本的仿射空间 $G/N$，看作 $G\times A$-簇（$A$ 是极大环面）。

\begin{thm}
  $X$ 是（拟仿射）球 $G$-簇，当且仅当 $F[X]$（$X$ 的正则函数环）作为 $G$-模是重数自由的。
\end{thm}
\begin{proof}
  （$\Rightarrow$。）固定 $B\subset G$。此后我们假定 $G$ 从右作用。注意 $F[X]$ 是 $G$ 的（无限维）局部有限表示，因此分解为最高权模 $V_\chi$ 的直和，其中 $\chi$ 是 $B$ 的特征标（即 $A$ 的特征标沿自然满射 $B\twoheadrightarrow A$ 拉回）。我们要证明 $V_\chi$ 的重数为 $1$。设 $f_\chi$ 是 $V_\chi$ 中的最高权向量，看作 $X$ 上的函数。由于 $B$ 经 $\chi$ 作用在 $f_\chi$ 上，故对一切 $x\in X$、$b\in B$ 有 $f_\chi(xb)=\chi(b)f_\chi(x)$。因此 $f_\chi$ 在 $X$ 中任一 $B$-轨道上的取值由它在一点处的取值决定。而按假设 $X$ 有稠的 $B$-轨道，故 $f_\chi$ 在这个稠轨道上（从而在整个 $X$ 上）只差一个标量就被唯一确定。换言之，$F[X]$ 中任何别的 $(B,\chi)$-本征向量都只能与 $f_\chi$ 相差一个标量，这就证明了 $(B,\chi)$-本征向量在 $F[X]$ 中的重数为 $1$，从而 $V_\chi$ 在 $F[X]$ 中的重数为 $1$，如所愿。
\end{proof}

要研究球簇，我们考察两件事：(1) 稠 $B$-轨道，记作 $\mathring X$；(2) 补集 $X\setminus\mathring X$。首先，若 $x\in\mathring X$ 任意，则 $\mathring X\simeq B_x\backslash B$。注意换一个 $x\in X$ 会得到 $B$ 的共轭子群。

现在考虑 $X$ 上有理函数空间 $F(X)$。由稠密性，它与 $\mathring X$ 上的有理函数空间相同。再考虑 $F(X)$ 中 $B$-本征函数的集合，记作 $F(X)^{(B)}$。注意这些函数在 $N$ 的（右）作用下不变。我们有：
\begin{equation}
  F(X)^{(B)}\simeq F(B_x\backslash B/N)^{(A)}.
  \label{eq:8-20}
\end{equation}
由此可以定义 $A_X$，即「$X$ 的 Cartan」，它是 $A$ 在 $\mathring X/N$ 上忠实作用的最大商。注意这与 $B$ 和 $x$ 的选取无关，因为 $B$ 与 $B_x$ 都只确定到共轭，从而在 $G$ 的「那个」Cartan（即看作典范子商的极大环面）中，它们在模 $N$ 的意义下典范地相同。于是有满射
\begin{equation}
  A\twoheadrightarrow A_X,
  \label{eq:8-21}
\end{equation}
它给出对偶环面之间的映射
\begin{equation}
  \check A_X\to\check A.
  \label{eq:8-22}
\end{equation}
注意 (8.22) 一般\textbf{不是}单射！它是单射，当且仅当 $\ker(A\twoheadrightarrow A_X)$ 连通。

下一讲我们会把 $\check A_X$ 扩张成一个约化群 $\check G_X$，即 $\check G$ 的一个子群，其极大环面为 $\check A_X$。这就是上面讨论的球簇所对应的对偶群。

让我们简单谈一下 (2)，即 $\mathring X$ 在 $X$ 中的补集。它是 $B$-不变除子的并。可以发现，$X$ 上为 $B$-本征函数的正则函数，恰是 $F(X)^{(B)}$ 中在这个补集的 $B$-除子上取正赋值的那些元素。这些 $B$-除子的数据，以及它们与 $F(X)^{(B)}$ 的（以赋值为配对的）配对，是给出球簇组合描述的关键。

\begin{exmp}
  设 $X=\mathrm{GL}_n$，$G=\mathrm{GL}_{n-1}\times\mathrm{GL}_n$ 分别从左、右作用。确定 $F[X]$ 作为 $G$-模的结构，与计算限制 $\pi^{\mathrm{GL}_n}|_{\mathrm{GL}_{n-1}}$ 的所谓「分支问题」密切相关。
\end{exmp}
